\section{Generator cores and bounded-support reconstruction}\label{sec:support}

The local topology space is infinite because any directed generator segment may carry an arbitrarily long ordered word of port-bearing tree vertices.  The finite computation becomes complete only after a structural theorem shows that bounded restrictions recover every such word.  We derive that theorem here from the level-2 and strong tree-child conditions.

\subsection{The finite core list}

For a connected blob $B$, its cyclomatic number is
\[
\beta(B)=|E(B)|-|V(B)|+1.
\]
Choose any rooted partner and delete one incoming edge at each reticulation, as in the construction of a displayed tree.  The restriction of the resulting undirected graph to $B$ is a forest.  If $r(B)$ is the number of reticulations in $B$, at most $r(B)$ edges of $B$ were deleted, whereas at least $\beta(B)$ edge deletions are required to turn a connected graph into a forest.  Hence
\begin{equation}\label{eq:cycle-rank-reticulation}
\beta(B)\le r(B).
\end{equation}
A level-2 blob therefore has cyclomatic number at most two.  Delete pendant cut-edge components but mark their attachment sites as ports, and suppress every unported ordinary vertex of degree two in the blob.

\begin{proposition}[Level-2 core theorem]\label{prop:core-theorem}
Every nontrivial reduced level-2 blob core is one of the following.
\begin{enumerate}[label=\textup{(\roman*)}]
\item A cycle core with one source event and one path-sink reticulation event.
\item A theta core: two branch vertices joined by three internally disjoint paths, with one source event and two reticulation events.  Up to branch reversal and path permutation, acyclicity, reachability, and binary degree constraints give exactly four directed theta cores.
\end{enumerate}
Every full ported blob is obtained uniquely from one such directed core by inserting ordered words of ordinary port-bearing tree vertices on its directed segments and adding the compulsory children of path-sink reticulations.
\end{proposition}

\begin{proof}
Equation~\eqref{eq:cycle-rank-reticulation} gives $\beta(B)\le2$.  If $\beta(B)=1$, biconnectivity forces $B$ to reduce to a cycle.  Suppose $\beta(B)=2$.  After suppressing degree-two path vertices, every remaining core vertex has degree three: biconnectivity gives minimum degree at least two, all degree-two vertices have been suppressed, and binary network degrees give maximum degree three.  Thus the kernel is cubic.  Writing $v$ and $e$ for its numbers of vertices and edges,
\[
e-v+1=2,
\qquad
2e=3v,
\]
so $v=2$ and $e=3$.  The three kernel edges expand to three internally disjoint paths between two poles, which is a theta graph.

For orientations, record along the three theta paths the source event $S$ and the one or two path-sink events $X$, allowing a reticulation to occupy a branch vertex.  Binary indegree/outdegree constraints fix each segment direction adjacent to an event; empty paths have two possible directions.  Reject directed cycles and orientations not reachable from $S$.  Exhaustive branch reversal and path permutation reduce the $24$ valid raw orientations to four core types.  This elementary enumeration is regenerated by \texttt{core\_enumerator.py}.  Re-inserting suppressed degree-two vertices produces ordered words on directed segments, and each such vertex carries exactly one bridge port in a binary network.
\end{proof}

\begin{figure}[p]
\centering
\begin{tikzpicture}[scale=.72, every node/.style={font=\scriptsize}]
% helper positions manually
% cycle
\begin{scope}[shift={(0,3.2)}]
\node[treev] (s) at (-1.5,0) {$S$}; \node[reticv] (x) at (1.5,0) {$X$};
\draw[reticedge] (s) .. controls (-.5,.8) and (.5,.8) .. (x);
\draw[reticedge] (s) .. controls (-.5,-.8) and (.5,-.8) .. (x);
\node at (0,1.3) {cycle core};
\end{scope}
% theta 0
\begin{scope}[shift={(4.6,3.2)}]
\node[treev] (u) at (-1.6,0) {$U$}; \node[treev] (v) at (1.6,0) {$V$};
\node[root] (s) at (0,1.15) {$S$}; \node[reticv] (x1) at (0,-1.1) {$X_1$}; \node[reticv] (x2) at (0,0) {$X_2$};
\draw[dir] (s)--(u); \draw[dir] (s)--(v);
\draw[reticedge] (u)--(x2); \draw[reticedge] (v)--(x2);
\draw[reticedge] (u)--(x1); \draw[reticedge] (v)--(x1);
\node at (0,1.75) {theta core 0};
\end{scope}
% theta 1 schematic
\begin{scope}[shift={(9.2,3.2)}]
\node[treev] (u) at (-1.6,0) {$U$}; \node[treev] (v) at (1.6,0) {$V$};
\node[root] (s) at (-.45,1.15) {$S$}; \node[reticv] (x1) at (.45,-1.1) {$X_1$}; \node[reticv] (x2) at (1.6,0) {$V$};
\draw[dir] (s)--(u); \draw[dir] (s)--(x2);
\draw[reticedge] (u)--(x2); \draw[reticedge] (u)--(x1); \draw[reticedge] (x2)--(x1);
\draw[ordinary] (u) .. controls (0,-.35) .. (x2);
\node at (0,1.75) {theta core 1};
\end{scope}
% theta 2
\begin{scope}[shift={(2.3,-.35)}]
\node[treev] (u) at (-1.6,0) {$U$}; \node[reticv] (v) at (1.6,0) {$V$};
\node[root] (s) at (0,1.15) {$S$}; \node[reticv] (x) at (0,-1.1) {$X$};
\draw[dir] (s)--(u); \draw[reticedge] (s)--(v);
\draw[reticedge] (u)--(v); \draw[reticedge] (u)--(x); \draw[reticedge] (v)--(x);
\draw[ordinary] (u) .. controls (0,-.35) .. (v);
\node at (0,1.75) {theta core 2};
\end{scope}
% core3 exact
\begin{scope}[shift={(7.4,-.35)}]
\node[treev] (u) at (-1.8,0) {$U$}; \node[reticv] (v) at (1.8,0) {$V$};
\node[root] (s) at (0,1.3) {$S$}; \node[reticv] (x) at (0,-1.25) {$X$};
\draw[reticedge] (u)--(v) node[midway,above] {$UV$};
\draw[dir] (s)--(u) node[midway,above left] {$SU$};
\draw[reticedge] (s)--(v) node[midway,above right] {$SV$};
\draw[reticedge] (u)--(x) node[midway,below left] {$UX$};
\draw[reticedge] (v)--(x) node[midway,below right] {$VX$};
\node at (0,1.9) {theta core 3};
\end{scope}
\node[align=center,text width=10.7cm] at (4.6,-3.05) {Ports arise by subdividing directed core segments with ordinary tree vertices carrying cut-edge components. The four theta cores are the complete orientation list after branch reversal and path permutation; the diagrams are schematic, while core 3 is labelled in the segment convention used in the seven-port proof.};
\end{tikzpicture}
\caption{The cycle core and four oriented theta cores. The exact machine encodings, rather than the drawings, define the finite atlas.}
\label{fig:cores}
\end{figure}


\subsection{Automatic exclusion of multi-triangle blobs}\label{sec:automatic-triangle}

The finite core list alone does not assume a bound on the number of triangles.  The next result shows that weak tree-child rootability supplies that bound automatically.

\begin{lemma}[A bridge cannot enter a reticulation]\label{lem:no-incoming-retic-bridge}
Let $D$ be a rooted acyclic phylogenetic network.  If an arc $x\to r$ enters a reticulation $r$, then its underlying undirected edge is not a bridge.
\end{lemma}

\begin{proof}
Let $y\to r$ be the other incoming arc.  If the undirected edge $xr$ were a bridge, removing it would place the root and $x$ in one component and $r,y$ in the other.  A directed root-to-$y$ path must cross the bridge through $x\to r$ and then continue from $r$ to $y$.  Together with $y\to r$, this creates a directed cycle, a contradiction.
\end{proof}

\begin{proof}[Proof of \cref{thm:automatic-triangle-bound}]
Let $N\in\TCw$ and choose a tree-child admissible rooting $D$.  By \cref{prop:core-theorem}, every nontrivial level-2 blob reduces to a cycle or theta core.  A cycle core has only one simple cycle, so consider a theta blob.  Let the three internally disjoint pole-to-pole paths have lengths
\[
1\le \ell_1\le\ell_2\le\ell_3.
\]
Its three simple cycles have lengths $\ell_i+\ell_j$.  Because the underlying binary network is simple, at most one theta path has length one.  If two distinct cycles are triangles, two pairwise sums equal three.  The only possibility is
\[
(\ell_1,\ell_2,\ell_3)=(1,2,2).
\]
Thus the blob is the graph $K_4$ minus one edge.  Write $u,v$ for its poles, $a,b$ for the internal vertices of the two length-two paths, and
\[
E(B)=\{uv,ua,av,ub,bv\}.
\]
The vertices $a,b$ each have one external bridge, while all three edges incident with either pole lie in the blob.

Here $\beta(B)=2$, so \eqref{eq:cycle-rank-reticulation} and the level-2 bound force exactly two reticulations in $B$.  If the two reticulations are adjacent, then either their common edge is directed from one reticulation to the other, giving a reticulation child, or the root subdivides that edge, giving the root two reticulation children.  Both alternatives violate tree-childness.

The only nonadjacent vertex pair is $\{a,b\}$.  Suppose therefore that $a$ and $b$ are the reticulations.  By \cref{lem:no-incoming-retic-bridge}, their external bridges are their outgoing edges; neither bridge can bring the root into the blob.  The binary root must consequently subdivide one of the five internal edges.  If it subdivides $uv$, then both poles have children $a$ and $b$.  If it subdivides an edge between a pole and one reticulation, the opposite pole is a parent of both $a$ and $b$.  In every case a tree vertex has two reticulation children, again contradicting tree-childness.  Hence no tree-child rooting exists, contrary to $N\in\TCw$.
\end{proof}

\begin{figure}[H]
\centering
\begin{tikzpicture}[scale=.82, every node/.style={font=\small}]
\begin{scope}[xshift=0cm]
  \node[treev] (u) at (0,1.35) {$u$};
  \node[treev] (v) at (3.0,1.35) {$v$};
  \node[treev] (a) at (1.1,2.50) {$a$};
  \node[treev] (b) at (1.1,.20) {$b$};
  \draw[ordinary] (u)--(v);
  \draw[ordinary] (u)--(a)--(v);
  \draw[ordinary] (u)--(b)--(v);
  \draw[bridge] (a)--++(0,0.60);
  \draw[bridge] (b)--++(0,-0.60);
  \node[align=center,text width=3.35cm,font=\scriptsize] at (1.5,-1.05)
    {two triangles force\\path lengths $(1,2,2)$};
\end{scope}
\draw[-{Stealth[length=2.4mm]},thick] (3.55,1.35)--(4.15,1.35);
\begin{scope}[xshift=4.6cm]
  \node[reticv] (u) at (0,1.35) {$u$};
  \node[reticv] (v) at (3.0,1.35) {$v$};
  \node[treev] (a) at (1.1,2.50) {$a$};
  \node[treev] (b) at (1.1,.20) {$b$};
  \draw[reticedge] (u)--(v);
  \draw[ordinary] (u)--(a)--(v);
  \draw[ordinary] (u)--(b)--(v);
  \node[align=center,text=retic,text width=3.35cm,font=\scriptsize] at (1.5,-1.05)
    {adjacent reticulations:\\a reticulation child};
\end{scope}
\begin{scope}[xshift=9.2cm]
  \node[treev] (u) at (0,1.35) {$u$};
  \node[treev] (v) at (3.0,1.35) {$v$};
  \node[reticv] (a) at (1.1,2.50) {$a$};
  \node[reticv] (b) at (1.1,.20) {$b$};
  \draw[ordinary] (u)--(v);
  \draw[reticedge] (u)--(a);
  \draw[reticedge] (v)--(a);
  \draw[reticedge] (u)--(b);
  \draw[reticedge] (v)--(b);
  \draw[bridge] (a)--++(0,0.60);
  \draw[bridge] (b)--++(0,-0.60);
  \node[align=center,text=retic,text width=3.55cm,font=\scriptsize] at (1.5,-1.05)
    {nonadjacent reticulations:\\one pole has two reticulation children};
\end{scope}
\end{tikzpicture}
\caption{Automatic exclusion of multiple triangles.  A simple theta with two triangle cycles must have path lengths $(1,2,2)$ and underlying graph $K_4$ minus the edge $ab$.  Adjacent reticulations force a reticulation child (or a root with two reticulation children).  The only nonadjacent pair is $a,b$; their external bridges must point away from them, and every internal root position leaves one pole with both reticulations as children.}
\label{fig:multitriangle-exclusion}
\end{figure}


\begin{remark}[Independent finite replay]\label{rem:multitriangle-replay}
The structural proof is independently checked by a primitive rooting census that imports none of the atlas code.  Across all external and internal root sites on $K_4$ minus an edge, it finds $25$ binary acyclic rooted presentations in seven graph-symmetry orbits and zero tree-child presentations.  A separately written C++ implementation obtains the same zero-survivor conclusion.  These computations audit the case split; they are not a substitute for the proof above.
\end{remark}

\subsection{Strong repairs and rigid supports}

A directed core may fail the tree-child condition before ports are inserted.  An inserted ordinary port vertex on a segment repairs its tail by providing a tree child.  For each core, a \emph{minimum strong repair} is a minimum set of segments whose subdivision makes the local rooted factor tree-child and prevents a reticulation child of a reticulation.

For the four theta cores, the minimum rigid-support sizes are
\begin{equation}\label{eq:support-sizes}
3,\quad4,\quad3,\quad3,
\end{equation}
where every path-sink child is included and one ordinary port is chosen on each segment of a contained minimum repair.  The cycle core has a two-port support: its sink child and one ordinary side anchor.

\begin{lemma}[Monotonicity of strong occupancy]\label{lem:occupancy}
If a set of occupied directed segments contains a minimum strong repair, then inserting additional port vertices on any segments preserves the local tree-child conditions.  Conversely, every full strongly tree-child expansion contains at least one minimum repair.
\end{lemma}

\begin{proof}
Subdividing a segment by an ordinary tree vertex changes no reticulation indegree and introduces a vertex with one continuation child and one port child.  It therefore cannot destroy a tree child already present.  The obligations in the unexpanded core are finite: every source/tree vertex lacking an ordinary child requires at least one of its outgoing segments to be occupied, and every reticulation whose core child is reticulate requires its unique outgoing segment to be occupied.  A full strong expansion satisfies all obligations; deleting occupied segments greedily while retaining satisfaction ends at a minimal repair.
\end{proof}

\begin{definition}[Rigid support]
A \emph{rigid support} of a full ported blob consists of all compulsory path-sink children together with one labelled port on each segment of one contained minimum strong repair.  We retain the labels of the corresponding descendant bridge components.
\end{definition}

The exact directed roles of a rigid support have trivial pointwise stabilizer except for the simultaneous side swap of the cycle core.  This prevents independent local automorphisms from making probe identifications inconsistent.

\subsection{One- and two-port probes}

Let $S$ be a fixed labelled rigid support.  If $p$ is an additional port, restrict the blob to $S\cup\{p\}$ and suppress unlabelled degree-two vertices.  The resulting support-plus-one topology records the directed segment containing $p$.  If $p,q$ lie on the same segment, the support-plus-two restriction to $S\cup\{p,q\}$ records their order; if they lie on different segments, their segment locations already distinguish them.

\begin{theorem}[Bounded-support reconstruction]\label{thm:bounded-support}
Every finite port-labelled strongly tree-child nonroot cycle or theta blob is determined, up to port-labelled core isomorphism and ordinary triangle redirection, by its restrictions to:
\begin{enumerate}[label=\textup{(\roman*)}]
\item one rigid support;
\item that support plus each single additional port;
\item that support plus each pair of additional ports.
\end{enumerate}
Consequently, theta blobs require at most six outgoing ports and cycle blobs at most four outgoing ports for structural reconstruction.  Including the distinguished incoming boundary gives at most seven and five boundary ports, respectively.
\end{theorem}

\begin{proof}
By \cref{lem:occupancy}, every full blob contains a rigid support of the stated size.  Its labelled roles identify the directed core and one anchor on every required segment.  A support-plus-one restriction determines on which core segment each extra label lies: after suppression, the label appears in the unique interval between the same pair of core/support reference points.  For two labels on one segment, the support-plus-two restriction distinguishes the two possible directed orders.  These pairwise comparisons form a transitive tournament because they arise from an actual directed word, and hence reconstruct the unique total order on each segment.  The cycle's only core automorphism swaps its two sides simultaneously; labelled anchors either break this symmetry or leave exactly that allowed swap.  Theta automorphisms are accounted for in the canonical core encoding and, pointwise on a labelled support, are trivial.
\end{proof}

\subsection{Weak restrictions and completion}

A support selected in a source blob need not remain strong when the same descendant blocks are marginalized in a competing target.  Omitted path-sink children and omitted repair ports become unobserved dummy components.  Completeness therefore requires a weak-target completion theorem.

\begin{lemma}[Completion of selected restrictions]\label{lem:weak-completion}
Every selected restriction induced from a full strongly tree-child cycle or theta blob is obtained by:
\begin{enumerate}[label=\textup{(\alph*)}]
\item assigning selected labels to arbitrary directed core segments and an arbitrary subset of path-sink children;
\item adding an unobserved child at every omitted path sink; and
\item adding unobserved port vertices on a minimum set of segments that repairs the full topology.
\end{enumerate}
After marginalization, different choices of dummy vertices on the same segment affect the JC tensor only through products of their edge multipliers.
\end{lemma}

\begin{proof}
Suppress all unselected port-bearing vertices in the full target.  What remains is the directed core, the selected labels, and possibly gaps at compulsory sink children or tree-child obligations.  Restoring one dummy leaf at every omitted sink and one ordinary dummy port on each missing repair segment produces a full strong extension.  Conversely, deleting the dummy leaves recovers the selected restriction.  In the Fourier parameterization, consecutive ordinary edges with identical displayed-tree descendant-mask vectors occur in every monomial together, so their parameters enter only through their product.  The product map $(0,1)^k\to(0,1)$ is surjective.
\end{proof}

For most selected restrictions, a completion needs at most two new outgoing labels.  Core~3 has a genuine three-label case.  The completion census is
\begin{center}
\begin{tabular}{c@{\qquad}r}
\toprule
additional rigid-support labels & labelled presentations\\
\midrule
1 & $2{,}256$\\
2 & $1{,}920$\\
3 & $192$\\
\bottomrule
\end{tabular}
\end{center}
Thus seven outgoing ports are both sufficient by \cref{thm:bounded-support} and necessary for the finite weak-target classification.  The $192$ residual records are classified in \cref{sec:seven-port}.
